\section{Introduction}
Single-particle cryo-electron microscopy (cryo-EM) now routinely produces three-dimensional density maps of macromolecular assemblies at near-atomic resolution, but the biological interpretation of those maps hinges on fitting atomic models into the observed density and refining them against it. The difficulty of that step has grown rather than diminished with improved reconstructions: larger complexes are imaged in multiple conformations, local resolution varies across a single map, and composite maps are increasingly deposited in which different regions were reconstructed with different particle subsets and orientations. In this regime, a single rigid or locally blurred model is a blunt tool for probing the map, because it must simultaneously describe both tightly ordered cores and flexible peripheries.

A diverse set of density-guided refinement procedures has emerged to address this. Several treat the map as a scalar field from which a gradient provides a driving force towards higher density, exemplified by the flexible fitting of atomic structures with molecular dynamics pioneered in MDFF~\cite{trabuco2008} and by the Flex-EM framework built on restrained refinement and comparative modelling~\cite{topf2008}. Others optimise the correlation coefficient between map and model directly, as in the real-space refinement module of \textsc{Phenix}~\cite{afonine2018,liebschner2019}, or adopt Bayesian formulations that weigh the experimental map against prior stereochemical knowledge~\cite{bonomi2018}. Sampling strategies include molecular dynamics with explicit or implicit solvent, gradient minimisation, and normal-mode analysis with elastic-network models~\cite{tama2004}. Interactive model-building tools such as \textsc{Coot}~\cite{emsley2010} and \textsc{ISOLDE}~\cite{croll2018} remain central at high resolution, and visualisation packages such as \textsc{UCSF ChimeraX}~\cite{goddard2018,pettersen2021} couple fitting utilities to structure exploration. All of these methods, however, inherit a common assumption: that the relationship between map and model is adequately captured by a single structure blurred by a global or locally adaptive B-factor.

That assumption breaks wherever the underlying ensemble matters. Multibody refinement in \textsc{RELION}~\cite{nakane2018}, focused refinement, and composite-map deposition all produce maps whose local resolution reflects the motion of a part of the complex rather than the quality of the overall reconstruction. Side chains that visit multiple rotameric states, loops that refold between conformers, and nucleic-acid strands whose phosphate backbone samples a band of positions are systematically misrepresented by a single-structure fit. Ensemble refinement in X-ray crystallography~\cite{levin2007,burnley2012} showed decades ago that multi-conformer descriptions capture features single models cannot, and this has been a long-standing assumption in the NMR community. Cryo-EM is now at the resolution range where those lessons apply directly.

Here we propose TEMPy-REFF (REsponsibility-based Flexible-Fitting), a refinement scheme that treats the map as a mixture of one Gaussian per atom together with a uniform background component, and iterates an expectation-maximisation update that yields self-consistent estimates of atomic positions, per-atom B-factors, and a scalar background level. The formulation is Bayesian at each step but computationally light: position updates are intensity-weighted centres of mass over voxels assigned to each atom by its responsibility, B-factor updates are variances of the same weighted distributions, and the background parameter is the mean residual intensity. We couple the density-side updates to a forcefield (AMBER with the GB-Neck2 implicit solvent model~\cite{nguyen2013}) evaluated in \textsc{OpenMM}~\cite{eastman2017}, with CHARMM36~\cite{huang2013} available as an alternative. Because the B-factor is estimated directly from the map rather than from the model, we can perturb atoms according to their local spread, locally minimise each perturbed structure, and generate an ensemble whose averaged blurred map is a better descriptor of the experimental density than any single member.

The same responsibility calculation supports two tasks that have never had a principled cryo-EM solution. Per-chain responsibilities partition each voxel's intensity among chains, giving a segmentation in which no voxel is assigned arbitrarily to one chain and no seams appear at chain boundaries. For composite-map construction, we apply the inverse: multiple focused maps are combined with responsibility-weighted intensities, yielding a smooth transition at boundaries without the seams of simple addition or the dilution of averaging. We benchmark the method on 366 cryo-EM maps from EMDB at resolutions between 1.8 and 7.1~\AA, using the CERES~\cite{liebschner2019} re-refined dataset as a reference; we show case studies on rotavirus VP6, glutamate dehydrogenase, yeast RNA polymerase~III, a nucleosome--CHD4 complex, and SARS-CoV-2 RNA polymerase initialised from an AlphaFold2 prediction~\cite{jumper2021,baek2021}. Quality of fit is assessed against SMOCf~\cite{farabella2015,joseph2017}, Q-score~\cite{pintilie2020}, EMRinger~\cite{barad2015}, MolProbity~\cite{williams2018,chen2010}, and CaBLAM~\cite{williams2018}, and against the local resolution estimated by ResMap~\cite{kucukelbir2014}, with gold-standard FSC conventions~\cite{rosenthal2003}. Across the benchmark the ensemble representation achieves CCC close to 0.9 and remains near-constant across resolution bands, delivering an average improvement of 0.197 over deposited models (paired $t$-test $p=0.0006$) and 0.150 over CERES ($p=0.003$).

\section{Related Work}
Our method sits at the intersection of three lines of work: density-guided refinement, ensemble representation of macromolecular structure, and Gaussian mixture modelling of cryo-EM maps. In density-guided refinement, the dominant paradigm has been to treat the map as a scoring target and sample the atomic coordinates against it. MDFF~\cite{trabuco2008} and its derivatives attach a fictitious force proportional to the local density gradient; Flex-EM~\cite{topf2008,joseph2017} blends restrained refinement with comparative modelling; \textsc{Phenix} real-space refinement~\cite{afonine2018,liebschner2019} optimises a CCC-derived target with geometric restraints; \textsc{Coot}~\cite{emsley2010} and \textsc{ISOLDE}~\cite{croll2018} provide interactive regimes with immediate chemical feedback; and Bayesian methods~\cite{bonomi2018} combine the map with prior information in a single posterior. Ensemble refinement for macromolecular data has a long history outside cryo-EM: multi-conformer refinement in X-ray crystallography~\cite{levin2007,burnley2012} and population-of-models approaches in NMR have shown that a single structure cannot capture the variability hidden in low-resolution or dynamic regions, and similar observations are now emerging in cryo-EM. The crystallographic ensemble-refinement experience~\cite{burnley2012} also informs our plateau-based stopping criterion for ensemble size. Gaussian mixture modelling has been used to represent maps at coarse resolution~\cite{kawabata2008} and to carry out fitting at the level of a few hundred Gaussians per complex; our contribution is to push this representation to one Gaussian per atom and use it to estimate per-atom B-factors self-consistently. The analysis of local map quality relies on per-residue scores such as SMOCf~\cite{farabella2015,joseph2017}, Q-scores~\cite{pintilie2020}, EMRinger~\cite{barad2015}, and local resolution from ResMap~\cite{kucukelbir2014}; model-side validation uses MolProbity~\cite{williams2018,chen2010} and CaBLAM~\cite{williams2018}. Map-processing support comes from RELION multibody refinement~\cite{nakane2018,scheres2012} and the CCP-EM software suite~\cite{burnley2017,joseph2020}. Finally, recent advances in deep-learning-based structure prediction~\cite{jumper2021,baek2021} have changed what an ``initial model'' means: we show below that TEMPy-REFF integrates cleanly with AlphaFold2 predictions and can place confidently-predicted residues that were absent from the deposited PDB entry.
